\documentclass[aps,prd,twocolumn,superscriptaddress,nofootinbib]{revtex4-2}
\usepackage{amsmath,amssymb}
\usepackage{graphicx}
\usepackage{hyperref}
\usepackage{booktabs}
\usepackage{orcidlink}

\begin{document}

\title{The Dimension of Space from the Granularity of Hilbert Space}

\author{Andrew Korytko\,\orcidlink{0009-0005-4569-2228}}
\email{andrew@alphalatitude.com}
\affiliation{AlphaLatitude Inc., P.O. Box 64341, Sunnyvale, California 94088, USA}

\date{\today; DOI: \href{https://doi.org/10.5281/zenodo.22678537}{10.5281/zenodo.22678537}}

\begin{abstract}
We derive the three dimensions of space from two postulates. The first is the register: every qubit is a string of $L$ bits around a ring of $L$ cells of one Planck length, with $L = 6.4\times10^{61}$ universal. Its bits are the outcomes of one measurement and its position is the phase, and each cell counts phase from its own zero mark, on which no law depends. The second is that there is a sphere in $d$-dimensional space, $d$ unknown, every great circle of which is such a ring, and its cells label the momenta an elementary interaction can carry, one cell per momentum. An elementary interaction is an event at which its legs exchange momentum. Momentum is conserved, so the momenta of an interaction with $k$ legs span $k-1$ dimensions, and each of them is one of $L^{k-1}$. The sphere has $L^{d-1}$ cells, so $d = k$ for each interaction, and $d$ need not be universal. In the register a moving string can pass one step of phase to the offset of two neighboring cells. That event has three legs, so $d = 3$. The offsets are light. An observation is an exchange of momentum, and a string's momentum changes in no other way, except when two strings meet at one cell, which for uncorrelated strings happens once in $L$ pairs. Every observation is therefore made with this event. A four-leg interaction, rare or common, would be seen only through it, and its four momenta span three dimensions, so it would appear three-dimensional. With $L$ from the cosmological constant the sphere is the horizon of an observer at the interaction. A count relaxed to real $d$ finds no other solution.
\end{abstract}

\keywords{Rational quantum mechanics; discrete Hilbert space; dimension of space; de Sitter horizon; holography}

\maketitle

\section{Introduction}
\label{sec:intro}

The dimension of space is an input to every physical theory. Ehrenfest asked in 1917 what in the laws of physics makes it manifest that space has three dimensions~\cite{ehrenfest1917}. In the century since, the question has been answered in two ways. One is consistency: string theory fixes the dimension of its spacetime by requiring the theory to be consistent on an assumed background~\cite{lovelace1971,goddard1972}. In the hypergraph-rewriting models of Ref.~\cite{wolfram2020} it is left open. There, dimension is a property of the graph that emerges on large scales. It need not be an integer, and three is the value that a rule not yet found would have to produce.

The other way starts from the qubit, whose state space, the Bloch ball, is three-dimensional. M\"uller and Masanes derive $d = 3$, and the qubit with it, from operational postulates on a system that carries direction information and evolves continuously and reversibly~\cite{muller2013}. Masanes, M\"uller, P\'erez-Garc\'ia, and Augusiak derive the three dimensions of the Bloch ball from entanglement~\cite{masanes2014}. Daki\'c and Brukner derive $d = 3$ from the classical limit of the theory, for pairwise interactions, and note that interactions among more particles might open higher dimensions~\cite{dakic2016}. On a discrete substrate, D'Ariano and Perinotti obtain the Dirac equation in three space dimensions from the minimal quantum cellular automaton that is unitary, local, homogeneous, and isotropic~\cite{dariano2014}. The reconstructions assume that physical space is locally Euclidean, of unknown dimension $d$; the automaton is placed on a lattice in three. This paper assumes less than either: a linear structure at the point of an interaction, in which a direction is fixed by $d-1$ numbers and vectors add, and no manifold (Sec.~\ref{sec:discussion}); whether even that linearity can be derived is asked in Ref.~\cite{muller2017}. This paper belongs with the second way, with the continuum given up. The qubit is taken from quantum mechanics, the input is a discrete register, and the dimension is a count, from two postulates. The loophole that Daki\'c and Brukner leave open is what the count measures: an interaction with $k$ legs needs $k$ dimensions (Sec.~\ref{sec:vertex}), and every observation is made with the one that has three.

The framework has two sources, and Sec.~\ref{sec:postulates} defines the terms used here. Palmer's Rational Quantum Mechanics (RaQM)~\cite{palmer2020,palmer2026,palmer2026b,palmer2026c} admits a qubit only where its two coordinates on the Bloch sphere, the sphere of a qubit's states, are rational with denominator $L$, and writes it as a string of $L$ bits. The register, Postulate~1 below, makes $L$ universal and puts the string on a ring of $L$ cells, each a Planck length long, on which turning the ring by one cell is a translation by $\ell_P$; Born's rule, the rationality condition, and the momentum quantum follow from it (Secs.~\ref{sec:postulates} and \ref{sec:register}). A companion paper~\cite{korytko2026c} develops the dynamics along one ring from the same postulate, and nothing here depends on it. The count below does not use the value of $L$, which enters only the numbers quoted, the root of Sec.~\ref{sec:root}, and the identification of Sec.~\ref{sec:pi}. It is fixed by placing the ring on a great circle of the cosmological horizon, one Planck length per step of phase,
\begin{equation}
R_\Lambda = \frac{L\,\ell_P}{2\pi},
\label{eq:ident}
\end{equation}
with $R_\Lambda = \sqrt{3/\Lambda}$, which gives $L = 6.4\times10^{61}$ for the observed $\Lambda$~\cite{planck2018}. A register is a statement about a circle. It says nothing about how many rings there are or how they fit together, and it leaves the three-dimensional bulk unconstructed.

The first postulate of this paper is the register, with each cell carrying its own zero mark for phase, on which no law depends. The second is about the sphere the circles make up. There is a sphere in $d$-dimensional space, with $d$ unknown, every great circle of which is such a ring, and its cells label the momenta that an elementary interaction can carry, one cell per momentum. An elementary interaction is a single event at which momentum is exchanged, and the carriers of momentum that enter and leave it are its legs. The least interaction that the first postulate allows has three legs. As a string moves from one cell to the next, the phase that arrives is counted from the next cell's mark, so the move adds to the phase an amount that belongs to the pair of cells, their offset. Where nothing else has happened, the offset is the difference between the two marks. In that interaction one notch of phase, one of the $L$ steps around the ring, passes between the string and the offset, and the offsets carry the momentum that the string gains or loses.

The derivation is a count. Momentum is conserved, because no law depends on the zero marks, and three momenta that sum to zero lie in a plane; Sec.~\ref{sec:vertex} shows that for this interaction they do not lie on a line. A momentum in the plane has one of $L$ headings and one of $L$ magnitudes, $L^2$ in all, and a sphere in $d$ dimensions whose great circles are rings has $L^{d-1}$ cells, so one cell per momentum gives $d = 3$. The same count gives an interaction with $k$ legs the dimension $k$, so the dimension belongs to the interaction and is not assumed to be universal. Space is observed to be three-dimensional because every record is made with this event: the offsets are light, and in the register a string's momentum changes in no other way, except when two strings meet at one cell.

An interaction with four legs is not forbidden, and the postulates do not exclude regions where such interactions are common. Its momenta would need four dimensions, but it could be seen only from outside, through the three-leg interaction with which every record is made, and its four momenta, which obey one linear relation, span three dimensions, so it would appear three-dimensional. In the register it needs two strings at one cell, which for uncorrelated strings happens for one pair in $L$, and the rational grid has exceptions of the same size, the eight angles of Niven's theorem among the $L$ angles of a ring.

Can some other $d$, integer or not, do as well as three? The question is not whether a world of non-integer dimension is possible, which the postulates cannot address, but whether the counting equation has a second solution once $d$ is allowed to be real. Section~\ref{sec:root} relaxes the rationality condition on purpose and counts the sphere in pixels of one Planck area, a count defined for real $d$. The exact count gives three at the finite $L$; the relaxed equation has one solution, $3.008$, and the $0.008$ is the known factor $\pi$ between the two counts, expressed in units of a dimension, so the relaxation finds no other dimension.

Section~\ref{sec:postulates} states the setting, the definitions, and the two postulates, and Sec.~\ref{sec:register} lists the results that one ring gives. Section~\ref{sec:vertex} derives conservation, the dimension of an interaction, and the register's three-leg interaction. Sections~\ref{sec:supply} and \ref{sec:root} count the cells of a sphere and relax the count to real $d$. Sections~\ref{sec:pi} and \ref{sec:states} identify the sphere with the horizon and show why a rival number, $L^3$, does not select $d = 4$. Sections~\ref{sec:algebra} and \ref{sec:consequences} give an algebraic reading and the consequences for the register, Sec.~\ref{sec:pred} the falsifiers, and Sec.~\ref{sec:discussion} closes with a ledger of what is postulated, derived, and open. Three things the paper does not do are stated now so that they are not read into it. It does not construct space or assume a manifold; it does not give the offsets a dynamics; and it does not claim that a physical principle excludes non-integer dimensions. The ledger lists its limits in full.

\section{The Setting and the Two Postulates}
\label{sec:postulates}

\textit{The setting.} The question is asked on an assumed background, as the critical dimension of string theory is~\cite{lovelace1971,goddard1972}. Only the structure of space at the point where an interaction takes place is used. A direction at that point is fixed by $d-1$ numbers, and a quantity with a magnitude and a direction there is a vector with $d$ components. Vectors at a point add, and that linear structure is assumed as well. The number $d$ is the number of components of such a vector. The derivation uses that linear structure, and the geometry of the sphere of Postulate~2, which the postulate fixes, and no other property of the background. Whether space away from the point is flat or curved does not enter. Every translation in this paper is by a whole cell of the register, and no continuum limit is taken.

The number $d$ is not assumed to be the same at every interaction. Postulate~2 below gives a sphere to an elementary interaction and not to space, so what the postulates fix, interaction by interaction, is the set of momenta an interaction can draw on, and Sec.~\ref{sec:vertex} reads $d$ from that set. No manifold is used, so this is not the claim that a manifold changes its dimension from place to place. Every observation is made with one kind of interaction (Sec.~\ref{sec:vertex}), so observed space has the one dimension found for it. The paper does not construct the background. It asks which $d$ is consistent with the two postulates.

\textit{The qubit.} Quantum mechanics is taken with its Hilbert space. A qubit is a quantum system for which a measurement has two possible outcomes, written $1$ and $0$, or $+1$ and $-1$ when their mean is taken. The state of a qubit is a point $(\theta,\phi)$ of a unit sphere, the Bloch sphere, on which the measurement fixes an axis whose two poles are the two outcomes. The unit vector to the state's point is the Bloch vector. Its component along the axis, $\cos\theta$, is the mean outcome, and its azimuth $\phi$ about the axis is the qubit's phase. The same state is the ket
\begin{equation}
|\psi\rangle = \cos\tfrac{\theta}{2}\,|1\rangle + e^{i\phi}\sin\tfrac{\theta}{2}\,|0\rangle ,
\label{eq:qubit}
\end{equation}
whose two coefficients are its amplitudes. The ket is fixed up to an overall factor $e^{i\gamma}$, and $\gamma$ is the global phase. A phase is an angle, and an angle is counted from a zero mark whose position is a choice.

\textit{Momentum and interactions.} Momentum is what generates translation. It has a magnitude and a heading, the direction in which it points, so it is a vector. An event is what happens at one place at one time. An interaction is an event at which momentum is exchanged. Its legs are the carriers of momentum that enter and leave it, each with a definite momentum. A leg that leaves one interaction and enters another is one carrier, with one momentum at both, since a mode keeps its momentum along its ring (Sec.~\ref{sec:register}). Quantum mechanics assigns an amplitude to an event, and the amplitude of a single event is called a coupling. An interaction is elementary if it is a single event, so that its amplitude is one coupling and not a product of couplings. An elementary interaction is also called a vertex.

\textbf{Postulate 1 (the register).} \emph{Quantum mechanics is granular. Every qubit is a string of $L$ bits, with $L\in\mathbb{N}$ universal, around a ring of $L$ cells, each one Planck length long, which light crosses in one Planck time. The bits are the possible outcomes of one measurement of the qubit, and the string's cyclic position on the ring is the qubit's phase. Turning the ring one notch, the shift, is a translation by one Planck length. Each cell has its own zero mark for phase, and no law depends on the marks.}

\textbf{Postulate 2 (the sphere).} \emph{There is a sphere in $d$-dimensional space, $d$ unknown, every great circle of which is such a ring, and its cells label the momenta an elementary interaction can carry, one cell per momentum.}

The companion paper states the same postulate with its last sentence narrowed to the global phase of a state; here the zero mark of phase may be set separately at every cell, and Sec.~\ref{sec:vertex} shows what that freedom yields. The postulate leaves out the placement of the ring on the cosmological horizon, to which Sec.~\ref{sec:pi} returns. The string is Palmer's bit string~\cite{palmer2026}, and the word has that meaning throughout. It does not record successive measurements: its bits are alternatives for one measurement, and its index is the cell on the ring. Let $m$ be the number of 1s in the string, and let $n$ be its cyclic position, the number of cells, modulo $L$, from an origin cell to the start of the string. For a qubit read alone the canonical form of the string is a single run of $m$ 1s, and $n$ is where the run starts. In the register the place of an event is a cell, and its time is one Planck time, a tick.

Probabilities need no postulate, because the bits are the possible outcomes of one measurement. The frequency of the outcome 1 over the string is $m/L$, and the mean outcome is $2m/L - 1$. The mean outcome is the component $\cos\theta$ of the Bloch vector, so $\cos\theta = 2m/L - 1$, and the half-angle identity gives $\cos^2(\theta/2) = m/L$. The frequency of an outcome therefore equals its squared amplitude in Eq.~(\ref{eq:qubit}). That is Born's rule, obtained here as in Palmer's construction~\cite{palmer2026}. The phase is the cyclic position, $\phi = 2\pi n/L$. A notch is the turn of the ring by one cell, and since $L$ of them make a full turn, it is also the step $2\pi/L$ of phase, the smallest there is: phase comes in whole notches. Together,
\begin{equation}
\cos^2\tfrac{\theta}{2} = \frac{m}{L},\qquad \frac{\phi}{2\pi} = \frac{n}{L},
\label{eq:raqm}
\end{equation}
which is Palmer's rationality condition, and the points it admits are the rational grid on the Bloch sphere. The postulates give the frequencies of outcomes and do not say which alternative a measurement returns.

The last sentence of the postulate concerns the zero from which a phase is counted. A string can be found at any cell of its ring (Sec.~\ref{sec:register}), and its phase as found at a cell is counted from that cell's zero mark. Moving every mark by the same amount changes the global phase of a state and nothing else, so no law depends on the global phase. In Palmer's construction the global phase is a permutation $\xi$ of the bits, specific to the system but unknown~\cite{palmer2026c}, and the postulate hides it in the same way. The sentence says more, because the marks may also be moved one cell at a time, and Sec.~\ref{sec:vertex} shows that the register's interaction comes from this freedom. Where the marks can all be set alike, the register is that of the companion paper. The arrangement of an individual string is its canonical form under that permutation, so the exact bits cannot be obtained, and a qubit's state is the pair $(m,n)$. A state of many qubits is many strings read jointly, cell by cell, in blocks (Sec.~\ref{sec:consequences}). The relative arrangement of the strings is the correlations, and it is data.

Most of Postulate~2 is geometry and Postulate~1. A sphere consists of great circles, one in every plane through its center, and the postulate makes each of them a ring of Postulate~1, with $L$ cells one Planck length long, which fixes the size of the sphere. Its circumference is $L\ell_P$ and its radius is $R = L\ell_P/2\pi$, which is Eq.~(\ref{eq:ident}) read as a definition of $R$. Every great circle has $L$ cells, and with Palmer's rule applied to each of its $d-1$ angles (Sec.~\ref{sec:supply}) the sphere has
\begin{equation}
N_{\rm cell}(d) = L^{\,d-1}
\label{eq:cells}
\end{equation}
cells. On a single ring Postulate~1 already gives one cell per momentum, $L$ cells for the $L$ momenta of Sec.~\ref{sec:register}. The $L$ cells and the $L$ momenta of a ring are the two bases of one space of $L$ dimensions, related by the discrete Fourier transform, so on a ring the correspondence is Fourier duality. A heading is a direction, which is a point of the sphere of directions, so a cell that labels a heading is a label of momentum and not of position. What Postulate~2 asserts is the same correspondence on the sphere, for the momenta of an elementary interaction counted with their magnitudes: a magnitude and a heading together take one cell. The correspondence runs both ways, one cell for each momentum the interaction can carry and one such momentum for each cell. It is between labels, and it is not a map of a momentum onto the cell that lies in its own direction: the momenta of an interaction fill the linear span of its legs, which need not be the whole of space, while the cells of the sphere cover every direction, and the postulate asserts that the two sets have the same size.

The postulate is stated for an interaction. It does not give the value of $d$; Sec.~\ref{sec:vertex} shows that it gives each interaction as many dimensions as it has legs, so the value three comes from the legs and not from the postulate. It does not say that two interactions share a sphere, so the dimension belongs to the interaction, and two with different numbers of legs cannot share one (Sec.~\ref{sec:vertex}). It does not say where the sphere is, and it is read here as a sphere about the place of the interaction, so that a heading, a direction from that place, is a point of it. Until Sec.~\ref{sec:pi} the sphere is used only through its cells and its great circles; there the value of $L$ makes it the horizon of an observer at the interaction.

\section{Results from One Ring}
\label{sec:register}

Five results about a single ring are used below. The first three follow from Postulate~1 in a few lines. The last two also use the Dirac Hamiltonian, which is taken from quantum mechanics. The companion paper develops all five along one line, and nothing below depends on it. Of these, the doublet enters the count at one point, the two energy branches in the remark on the factor of two at the end of Sec.~\ref{sec:vertex}, and it is the object of Sec.~\ref{sec:algebra}; the tick is used in Secs.~\ref{sec:vertex} and \ref{sec:consequences}.

\textit{Momentum along a ring takes $L$ values.} A particle on the ring is described by the states $|j\rangle$, one for each cell $j$, counted modulo $L$. The shift acts on them as $X|j\rangle = |j+1\rangle$, and by Postulate~1 it is a translation by $\ell_P$. Momentum is what generates translation, $X = e^{-i\hat p\ell_P/\hbar}$, so a state of definite momentum is an eigenstate of $X$. The ring closes after $L$ cells, so $X^L = 1$ and the eigenvalues are the $L$th roots of unity, $e^{2\pi i\kappa/L}$ with $\kappa$ a whole number in $-L/2 < \kappa \le L/2$. The eigenstates are the modes $|\tilde\kappa\rangle$, waves with $\kappa$ wavelengths around the ring, and their momenta are the $L$ values $p_\kappa = h\kappa/L\ell_P$, in quanta of $2\pi\hbar/L\ell_P$, which is $\hbar/R_\Lambda$ once the ring is on the horizon (Sec.~\ref{sec:pi}). One shift advances the phase of a mode by $\kappa$ notches, since the mode has $\kappa$ wavelengths around the ring and the shift moves it by one cell; the pattern travels one cell per tick, whatever $\kappa$ is. A massless mode is one bit string, the string of its momentum doublet (below), with every bit the mode's sense; for a massive mode the doublet's string exists only at the ticks at which its reading lies on the grid. A string in a mode is what is called below a string with a definite momentum. The shift changes a mode by a phase only, so a mode is found at every cell with the same probability, which is $1/L$.

\textit{One ring for each direction of motion.} Postulate~1 gives a ring wherever there is translation, and a direction of motion is a line of translation. Beyond one line, therefore, every direction of motion has its own ring, and a mode advances along its own direction. The translation along a direction is the shift of that direction's ring.

\textit{The orientation of the grid is a choice of coordinates.} Every quantity in Eq.~(\ref{eq:raqm}) is an overlap of two states or a phase relative to a reference. The orientation of the grid in space is therefore not data.

\textit{The momentum doublet.} A particle on the ring moves around it in one of two senses. It is described by the Dirac Hamiltonian in one dimension, $H = c\hat p\sigma_z + Mc^2\sigma_x$. The Pauli matrix $\sigma_z = \pm1$ labels motion to the right or to the left at the speed of light, and $\sigma_x$, with the mass $M$, reverses the sense. The two senses at a given momentum are the two outcomes of one qubit, the particle's chirality, which this paper calls the momentum doublet. Its number of 1s gives the frequency of moving right, $m/L$, so that $v/c = 2m/L - 1$. The energies are $\pm E$ with $E = \sqrt{p^2c^2 + M^2c^4}$, the two energy branches, and in a state of definite energy $\cos\theta = pc/E$ on the positive branch and $-pc/E$ on the negative one.

\textit{The tick.} Time advances in ticks of one Planck time, and each tick is the shift followed by the flip, $U = e^{-i\alpha\sigma_x}X^{\sigma_z}$. The shift moves every bit of a string one cell along the ring in its sense. The flip reverses a string's sense with amplitude $\sin\alpha$, where the rate $\alpha = M/m_P$ is the mass in units of the Planck mass, taken from the Dirac Hamiltonian with $M$ free; the companion paper obtains this form of the tick from Postulate~1 with the mass term. The flip moves no momentum between strings.

\section{The Dimension of an Interaction}
\label{sec:vertex}

Postulate~2 ties the cells of the sphere to the momenta of an interaction. This section counts those momenta, which needs the conservation of momentum first, and then reads off the dimension.

\textit{Conservation.} A phase is a cyclic position (Postulate~1), so the global phase of a state is the common origin of its strings on a ring. Moving every zero mark by the same amount changes that origin and nothing else, and by the last sentence of Postulate~1 no law depends on the marks. No law therefore singles out a first cell, and every law is unchanged when everything on a ring, the strings and the offsets defined below, is moved along together by one cell, which means that it commutes with the joint shift $X$. The eigenvalue of $X$ is the mode number $\kappa$, and for several carriers it is additive, so the total $\kappa$ is conserved. That is Noether's theorem on $\mathbb{Z}_L$: the translation invariance of the ring is the conservation of momentum along it. It holds modulo $L$: a total mode number that passes $L/2$ wraps around, which on a lattice is called umklapp, and the wrap adds a momentum along the ring. The wrap changes what a sum of momenta comes to and not how many momenta there are, so it leaves the counting below untouched. The ring gives the law along one direction. Every direction of motion has its own ring, and the component of a momentum along a direction is what generates the translation along that direction's ring (Sec.~\ref{sec:register}). The total component is therefore conserved along every direction, and the momenta at a vertex sum to zero as vectors.

\textit{The orientation of a vertex.} The orientation of the grid is a choice of coordinates (Sec.~\ref{sec:register}), so for a single vertex only the relations among its legs are data. The orientation of the span of its momenta is a choice, and a leg's momentum is described within the span, by a magnitude and a heading. For several vertices in several planes only the overall orientation is a choice; Sec.~\ref{sec:states} takes that case up.

\textit{A vertex with $k$ legs needs $k$ dimensions.} Let a vertex have $k$ legs, and let $s$ be the dimension of the space its momenta span. Within that span a momentum has a signed mode number, with $L$ values (Sec.~\ref{sec:register}), and a heading. The headings within the span are the points of the great sphere in which the span meets the sphere of Postulate~2. That great sphere has $s-1$ angles, each with $L$ values when the grid is aligned with the span (Sec.~\ref{sec:supply}), so the headings number $L^{s-1}$. The momenta a leg can carry therefore number $L^{s}$, and this set is called the alphabet of the vertex. By Postulate~2 the alphabet and the cells of the sphere correspond one to one, so $L^{s} = L^{d-1}$, and the dimension is one more than the span, $d = s+1$. The step from the alphabet to the cells is Postulate~2 itself, and nothing else licenses it. The $d$ found this way is the $d$ of the setting, the number of components of a vector at the point of the vertex; no metric and no wave equation in $d$ dimensions is claimed or used. The momenta of a vertex span a subspace one dimension smaller than its space, and never the whole of it.

The legs enter through the span. Momentum is conserved, so the $k$ momenta obey one linear relation and $s \le k-1$, with equality when they obey no other relation. In that case
\begin{equation}
d = k ,
\label{eq:nk}
\end{equation}
and a vertex needs as many dimensions as it has legs. A vertex whose momenta obey a further relation has a smaller span and a smaller $d$, and that case is settled for the register's interaction below. The correspondence is used in both directions. A sphere with more cells than the alphabet would have cells that label no momentum the vertex can carry, and a sphere with fewer could not label the alphabet, so $d$ can neither exceed $s+1$ nor fall below it. What the equation weighs is the magnitude. The headings and the cells are both angles on the sphere, and they cancel; the factor of $L$ that remains is the magnitude, which comes from the modes of a ring (Postulate~1) with no $d$ in it. A magnitude and a heading together take one cell, and that is why the span falls one dimension short.

\textit{The offsets.} The theorem leaves the number of legs open. Postulate~2 speaks of elementary interactions, so it takes for granted that they occur, and the last sentence of Postulate~1 says what the least of them is in the register. The shift moves a string from one cell to the next, and the phase that arrives is counted from the next cell's mark. The move therefore adds to the phase an amount that belongs to the pair of cells, which is called their offset. Where nothing else has happened, the offset is the difference between the two marks. That no law depends on the marks is a local gauge invariance, and the offsets are what such an invariance brings with it, the link variables of a lattice gauge field~\cite{peierls1933,wilson1974,kogut1975}. Phase comes in whole notches (Postulate~1), so an offset is a whole number of notches. If the mark of one cell is moved by $c$ notches, then on every ring through that cell the offset on one side rises by $c$ and the offset on the other side falls by $c$, and nothing else changes. Two facts follow from this rule. The sum of the offsets around a closed path of cells is the same for every choice of marks, because every cell on the path is entered once and left once. And no choice of marks changes one offset alone.

\textit{The register's interaction has three legs.} An interaction exchanges momentum, and a string's momentum is the phase that it gains per cell (Sec.~\ref{sec:register}). A string therefore takes part in an interaction when, in some move, it gains more or less phase than its momentum provides. Phase comes in whole notches, so the least such amount is one notch. A move involves two neighboring cells, so unless a second string is present the only other party to it is their offset. The least interaction is therefore this one: as a string moves from one cell to the next, one notch passes between its phase and the offset of the pair. The notch that passes is called a quantum. The change in the offset is not a relabeling: it alters by one notch the sum of the offsets around every closed path through that pair of cells, and those sums are the same for every choice of marks. The offsets are therefore a carrier in their own right, and conservation fixes only how much they carry, what the string gains or loses. The event has three carriers of momentum, the string before, the string after, and the offsets, and these are its three legs. The amplitude of the event, which is its coupling, is not fixed by the postulates. Field theory on a lattice has the same vertex, with the offsets as its gauge field.

\textit{The third leg needs the sphere.} The event above takes place on one ring, and the offsets of one ring cannot be a leg. A leg carries momentum from place to place, and the only closed path on a ring is the ring itself, so the only datum in its offsets is their total, and a total has no place. On the sphere of Postulate~2, in more than two dimensions, the rings cross one another. Two of them meet at two opposite cells and enclose nothing shorter than a ring, but three or more enclose a triangle, so there are closed paths made of arcs of three rings and much shorter than a ring. Each such path has its own sum, which no choice of marks changes and which belongs to a definite place on the sphere. These sums are the whole of what the offsets hold: a set of offsets whose sum vanishes around every closed path, the rings taken whole included, is a choice of marks and nothing else, and it carries no momentum. The leg is an excitation of the sums. A quantum changes the sums of the paths that run through its pair of cells, and the passing of such changes from place to place is a wave of the offsets. A wave that runs along a ring has a whole number of wavelengths around it, as a mode of a string has, so the third leg's momentum takes the same values as the other two, which is all that the count needs from it. The gauge field of a free choice of phase is the electromagnetic field (Ref.~\cite{weinberg1995}, Ch.~8), so these waves are identified with light, and that identification is used below only in the statement that observations are made with light. How the waves travel, what their equation of motion is, and what the coupling is worth are not derived here, and none of them enters the count. They belong to a paper on the field.

\textit{The three legs are not collinear.} The count uses the span, so the case in which the three momenta lie on one line has to be settled. Were the momenta this interaction can carry confined to one line, they would be the $L$ signed modes of one ring, and by Postulate~2 the sphere would have $L$ cells, which is $d = 2$. In two dimensions the sphere is a single great circle, so every cell of it lies on one ring and the only closed path of cells is that ring taken whole. By the paragraph above, the offsets would then hold one number, their total, which belongs to no place, and none of them could carry momentum. The interaction would have no third leg, against the three it was just shown to have. Its momenta therefore span a plane. No two of them are collinear either, since two collinear momenta and the sum rule would put the third on the same line. The string therefore arrives on one ring, leaves on a second, and the quantum departs on a third, all three through the cell at which the event takes place. Which leg takes how much is the dynamics, and the count does not need it.

\textit{Space has three dimensions.} For the register's interaction $k = 3$, and Eq.~(\ref{eq:nk}) gives
\begin{equation}
d = 3 .
\label{eq:three}
\end{equation}
Observation agrees (Sec.~\ref{sec:pred}). The dimension belongs to the interaction, and nothing here makes it universal.

\textit{Why every observation is three-dimensional.} A record, in an eye or in an apparatus, is a change in the recorder that depends on what it records. Apart from the offset a string crosses, the tick moves and flips each string by itself (Sec.~\ref{sec:register}), so a change in one string that depends on another needs an exchange between them, and an exchange is an interaction at a cell. At a move the parties to an exchange are the string, the offset of its pair of cells, and a second string if one is at the same cell (above). A string's momentum therefore changes in two ways only: it passes a quantum, which is the three-leg vertex, or it meets another string at its cell. For uncorrelated strings the meeting is a coincidence, and the next paragraph shows that it happens for one pair in $L$. Every record is therefore made with the three-leg vertex, up to one pair in $L$, and that vertex has $d = 3$. Since the quantum is light (above), this is the statement that observation is made with light. The postulates do not exclude a vertex with four legs and a coupling of its own, or regions where such vertices are common. Rare or common, it would be seen only from outside, through three-leg vertices on its legs, since a leg has one momentum at both of its ends (Sec.~\ref{sec:postulates}). Its four momenta obey one linear relation and span three dimensions, so an outside observer would record four vectors in three-dimensional space that sum to zero, and nothing four-dimensional would appear. The dimension of that vertex's own sphere could not be observed, and its only trace would be a coupling of its own.

\textit{Four legs in the register.} By Eq.~(\ref{eq:nk}) the momenta of a vertex with four legs would need four dimensions and a sphere of $L^3$ cells of its own; two vertices with different numbers of legs cannot share one sphere, since Eq.~(\ref{eq:nk}) would then have two values. In the register an event with four legs at one cell in one tick arises in two ways. One string may pass two quanta in one move, with amplitude the product of the two couplings; these are the four-point terms of field theory that a cubic coupling fixes. Or two strings may be at the same cell in the same tick, where each passes a quantum or the two exchange a notch directly. That is a coincidence, and its rate can be computed. A leg has a definite momentum, so it is a mode, and a mode is found at each cell with probability $1/L$ (Sec.~\ref{sec:register}). For two uncorrelated strings on one ring, the probability that both are found at the same cell is
\begin{equation}
\sum_{j=1}^{L}\frac{1}{L}\cdot\frac{1}{L} = \frac{1}{L} ,
\label{eq:coincidence}
\end{equation}
which is exact. The count concerns two events at one cell and not their order, so it is the same in both directions of time. Two rings in different planes share at most two cells, and for strings on them the probability is at most $2/L^2$. A coincidence is therefore at most one in $L = 6.4\times10^{61}$ of the pairs of events. If it is counted as one vertex with four legs, its momenta need four dimensions, and by the argument above it would still appear three-dimensional. When each string passes a quantum, the events are two whether they coincide or not, each with the alphabet of Eq.~(\ref{eq:demand}), and their momenta form a configuration in the sense of Sec.~\ref{sec:states}. The rule is per vertex and not per process: a $2\to3$ process is a chain of three-leg events in the same way.

The rational grid has exceptions of the same size. By Niven's theorem~\cite{niven1956,jahnel2010} an angle $\alpha$ has $\alpha/2\pi$ and $\cos\alpha$ both rational at eight angles only, the multiples of $\pi/2$ and of $\pi/3$. On a ring these are at most eight of the $L$ angles, a fraction $8/L$, and only at these can the angle between two measurement settings and its cosine both be rational, which is what Palmer's Impossible Triangle Corollary~\cite{palmer2026b,palmer2026c} turns on. In both cases the exceptions exist, they are of order one part in $L$, and nothing can detect them.

\textit{A check against field theory.} Field theory agrees that a record is made with a three-leg vertex. With fermions it has no four-point term that contains a photon, the quantum of light (Ref.~\cite{weinberg1995}, Ch.~8). Its four-point terms with two photons, the seagull terms, which occur for charged scalars and for the $W$, are fixed by the cubic coupling as its square, and in the register they are two events at one cell. The scattering of light by light is a loop of three-leg vertices. In field theory as well, then, every record is made with a three-leg vertex.

\textit{The alphabet of a three-leg vertex.} For the register's interaction the span is a plane, and the count reads label by label. The heading of a momentum in the plane is one of the $L$ cells of the great circle in which the plane meets the sphere, and its magnitude is one of the $L$ modes of the ring along that heading (Sec.~\ref{sec:register}), so
\begin{equation}
N_{\rm mom} = L\times L = L^2 = 4.10\times10^{123},
\label{eq:demand}
\end{equation}
in which $d$ does not appear. Were the orientation of the plane counted as well, with its $2(d-2)$ angles, the alphabet would be $L^{2d-2}$, which equals $L^{d-1}$ only at $d = 1$, where there is no plane. The count therefore depends on the orientation being a choice of coordinates, which it is for one vertex; for two vertices the relative orientation of their planes is data, and Sec.~\ref{sec:states} counts it separately.

\textit{No factor of two.} The heading is oriented and the mode is signed, and a reader might count headings over a half circle, giving $L^2/2$, on the ground that the two signs say the same thing twice. They do: a momentum with its sign and its heading both reversed is the same vector, so there are $L^2/2$ vectors. Each has two energy branches (Sec.~\ref{sec:register}), which are different states, and a vertex can carry both. The momenta of Postulate~2 are these states, and they number $L^2$ with no factor of two. The zero mode is the one place where the two counts part, since the $L$ headings at $\kappa = 0$ name one vector with two states and not $2L$; that is one label in $L$, the size of the exceptions above.

\section{The Cells of a Sphere in $d$ Dimensions}
\label{sec:supply}

Postulate~2 gives the sphere $L^{d-1}$ cells, Eq.~(\ref{eq:cells}). This section constructs those cells for any $d$ and counts the sphere a second way, in pixels of one Planck area. The grid gives the count that the postulates call for, and the pixel, which is defined below, continues that count to real $d$ in Sec.~\ref{sec:root}.

\textit{The rational grid.} Write $S^{d-1}$ in hyperspherical coordinates, $\theta_1,\dots,\theta_{d-2}\in[0,\pi]$ and $\phi\in[0,2\pi)$. The area element is
\begin{equation}
dA = R^{\,d-1}\prod_{i=1}^{d-2}\sin^{\,d-1-i}\theta_i\,d\theta_i\;d\phi,
\label{eq:dA}
\end{equation}
a product of factors each depending on one angle. For each polar angle let
\begin{equation}
u_i(\theta_i) = \frac{\int_0^{\theta_i}\sin^{d-1-i}\theta\,d\theta}{\int_0^{\pi}\sin^{d-1-i}\theta\,d\theta},\qquad u_{d-1} = \frac{\phi}{2\pi},
\label{eq:u}
\end{equation}
the fraction of the sphere's area within polar angle $\theta_i$ of the pole. Because the measure is a product, it is uniform in these variables: $dA = A_d\,du_1\cdots du_{d-1}$. On $S^2$ the $u$'s are $(1-\cos\theta)/2$ and $\phi/2\pi$, and these are Palmer's coordinates: $\cos^2(\theta/2) = 1 - u_1$ is the frequency $m/L$ of Eq.~(\ref{eq:raqm}), and $\phi/2\pi = n/L$. Palmer's rule is that every probability is a multiple of $1/L$, and each $u_i$ is a probability, that of a uniformly random direction lying within $\theta_i$ of the pole. Applied to each $u_i$, the rule gives $L$ values per coordinate and hence the product grid of Eq.~(\ref{eq:cells}), with $L^{d-1}$ cells, each of area $A_d/L^{d-1}$. On $S^2$ the $L$ values of $u_1$ cut the sphere into bands of latitude and those of $u_2$ into sectors of longitude, and a cell is one band crossed with one sector. At $d=3$ the grid is that of Eq.~(\ref{eq:raqm}), and its cell has area $4\pi R^2/L^2 = \ell_P^2/\pi$, the number that Sec.~\ref{sec:pi} explains. The grid is the natural extension of Palmer's rule to $d-1$ angles, and Sec.~\ref{sec:root} shows that the result does not depend on its details. The construction uses a whole number of angles. We have no version of it for non-integer $d$, and do not claim that none exists.

\textit{The grid is built on a pole.} Two partitions of the sphere are in play and should not be confused. Postulate~2 cuts every great circle into $L$ cells of one Planck length, which is what makes it a ring. The grid cuts the surface into $L^{d-1}$ patches of equal area, $L$ values in each of the $d-1$ coordinates. The two are not the same cut: at the equator a patch spans one Planck length along the circle and $1/\pi$ of one across it, and Sec.~\ref{sec:pi} gives the relation at $d = 3$. The grid is built on a pole, so the patches depend on that choice while their number does not, and the orientation of the grid is a choice of coordinates (Sec.~\ref{sec:register}). For a vertex the grid is turned so that the span of its momenta meets the sphere where the first $d-s$ polar angles are right angles; the $s-1$ angles that remain run over that great sphere, each with $L$ values. At $d = 3$ the span is a plane and it meets the sphere in the equator. The alignment costs nothing, because the orientation of a vertex is a choice as well (Sec.~\ref{sec:vertex}).

\textit{Planck pixels.} A pixel is a cell of one Planck area, $\ell_P^{\,d-1}$ in $d$ dimensions. It is a unit of area and not a cell of the grid, and at $d = 3$ it is the cell of the Bekenstein--Hawking entropy $A/4\ell_P^2$. The sphere's area is $\Omega_{d-1}R^{\,d-1}$, where $\Omega_{d-1}$ is the area of the unit sphere,
\begin{equation}
\Omega_{d-1} = \frac{2\pi^{d/2}}{\Gamma(d/2)},
\label{eq:omega}
\end{equation}
equal to $4\pi$ at $d=3$ and $2\pi$ at $d=2$, so
\begin{equation}
N_{\rm pix}(d) = \Omega_{d-1}\Big(\frac{R}{\ell_P}\Big)^{d-1} = \Omega_{d-1}\Big(\frac{L}{2\pi}\Big)^{d-1}.
\label{eq:supply}
\end{equation}
At $d=3$ this is $L^2/\pi = 1.30\times10^{123}$, the Planck-pixel count of the horizon, and at $d=2$ it is $L$.

\begin{table}[t]
\caption{Cells of the sphere against the $L^2$ momenta of a three-leg vertex, Eq.~(\ref{eq:demand}), at $L = 6.4\times10^{61}$.}
\label{tab:count}
\begin{ruledtabular}
\begin{tabular}{ccccc}
$d$ & $N_{\rm cell}(d)$ & $N_{\rm pix}(d)$ & value of $N_{\rm pix}$ & $L^2/N_{\rm pix}$ \\
\hline
2 & $L$ & $L$ & $6.4\times10^{61}$ & $6.4\times10^{61}$ \\
3 & $L^2$ & $L^2/\pi$ & $1.30\times10^{123}$ & $3.14$ \\
4 & $L^3$ & $L^3/4\pi$ & $2.09\times10^{184}$ & $2.0\times10^{-61}$ \\
5 & $L^4$ & $L^4/6\pi^2$ & $2.83\times10^{245}$ & $1.4\times10^{-122}$ \\
\end{tabular}
\end{ruledtabular}
\end{table}

\section{Non-Integer $d$}
\label{sec:root}

\textit{The question and the price of asking it.} On the rational grid the count reads $L^{d-1} = L^2$ for a three-leg vertex, and $d = 3$, Eq.~(\ref{eq:three}). The grid gives $L$ values to each angle, so $N_{\rm cell}(d)$ is defined at whole $d$ and nowhere else, and the question whether some other $d$, integer or not, would do as well cannot be put on the grid at all. This section puts it by relaxing the rationality condition on purpose. It tiles the sphere with the Planck pixels of Eq.~(\ref{eq:supply}), a count that is defined for real $d$, and looks for every solution of the relaxed equation. The relaxation has a price: the two counts differ by a factor that geometry fixes, known before any solution is computed, and the factor bounds how far a solution can lie from three.

Set the pixels against the $L^2$ momenta of the vertex,
\begin{equation}
\Omega_{d-1}\Big(\frac{L}{2\pi}\Big)^{d-1} = L^2 ,
\label{eq:match}
\end{equation}
first for integer $d$ (Table~\ref{tab:count}). Two dimensions fall short by a factor $L$, and four have $L/4\pi = 5.1\times10^{60}$ pixels to spare. Three is the only row in which the ratio is a number of order one and not a power of $L$.

\textit{The factor of the relaxation.} The pixel count and the grid count of Sec.~\ref{sec:supply} share the exponent and differ by the factor
\begin{equation}
C(d) = \frac{N_{\rm pix}(d)}{N_{\rm cell}(d)} = \frac{\Omega_{d-1}}{(2\pi)^{d-1}} .
\label{eq:C}
\end{equation}
Geometry alone fixes it. Postulate~2 sets the great circle at $L$ Planck lengths, so the area of the sphere is exactly $\Omega_{d-1}(L\ell_P/2\pi)^{d-1}$. Its values are $C(2) = 1$, $C(3) = 1/\pi$, and $C(4) = 1/4\pi$. Its logarithmic derivative is $\tfrac12\ln\pi - \tfrac12\psi(d/2) - \ln 2\pi$, where $\psi$ is the digamma function, and it runs from $-0.98$ at $d = 2$ to $-1.48$ at $d = 4$. So $C$ decreases on that interval and lies between $1/4\pi$ and $1$. Continued through its exponent alone, with $C$ set to one, the equation reads $d - 3 = 0$ at every $L$. Whatever departure from three appears below is produced by the pixel and by nothing else.

\textit{One root, within the bound.} Nothing in Eq.~(\ref{eq:match}) requires $d$ to be an integer, because the Gamma function is defined for real argument. With Eq.~(\ref{eq:C}) the equation reads $C(d)\,L^{d-1} = L^2$, and in logarithms
\begin{equation}
g(d) \equiv (d - 3) + \frac{\ln C(d)}{\ln L} = 0 .
\label{eq:g}
\end{equation}
The function $g$ is continuous, with $g(2) = -1$ and $g(4) = 1 - \ln(4\pi)/\ln L$, which is positive for every $L$ above $4\pi$. Its derivative, $1 + (\ln C)'/\ln L$, is positive on the interval for every such $L$. Equation~(\ref{eq:g}) therefore has exactly one solution $d_*$ between $2$ and $4$, which is called the root. At the root Eq.~(\ref{eq:g}) gives
\begin{equation}
d_* - 3 = \frac{\ln[1/C(d_*)]}{\ln L} ,
\label{eq:root}
\end{equation}
and because $1/C$ lies between $1$ and $4\pi$ on the interval,
\begin{equation}
0 \le d_* - 3 \le \frac{\ln 4\pi}{\ln L} = 0.018
\label{eq:bound}
\end{equation}
at $L = 6.4\times10^{61}$. The bound follows from the geometry of the sphere before the root is computed, so the argument is not circular. The left side of Eq.~(\ref{eq:match}) increases with $d$ from $d = 1$ until the growth of $\Gamma(d/2)$ overtakes it, near $d \sim 10^{123}$, so the root is the only solution below that value.

\textit{The value of the root.} The numerical root is $d_* = 3.008117$. At the root the ratio $1/C$ is $3.1745$, and Eq.~(\ref{eq:root}) returns $\ln(3.1745)/\ln L = 0.008117$, the same number. At $d = 3$ the ratio is exactly $\pi$, which gives the leading term $d_* \approx 3 + \ln\pi/\ln L = 3.0080$. The departure falls as $1/\ln L$, so the value of $L$ hardly matters: a factor of ten either way moves the root by less than $2\times10^{-4}$.

The answer to the question is therefore no. The relaxed equation has one solution, it lies inside the bound that the relaxation itself imposes, Eq.~(\ref{eq:bound}), and the relaxation turns up nothing that the grid had missed. No meaning is given to $d_*$. It is not a dimension and not a second value of $d$; it is the factor $\pi$ between a rational cell and a Planck pixel (Sec.~\ref{sec:pi}), expressed in units of a dimension. Two remarks complete the argument.

\textit{The choice of pixel cannot move the integer.} Another pixel, such as the $4\ell_P^2$ that holds one unit of horizon entropy, divides the count by a constant $c$ of order one and moves the root by $\ln c/\ln L$, which is $0.010$ for $c = 4$. To bring the root to a half-integer would take a factor $2.5\times10^{30}$.

\textit{The continuation is formal and standard.} Equation~(\ref{eq:omega}) continued through $\Gamma$ is the sphere area of dimensional regularization~\cite{thooft1972}, and the configuration is the same: external momenta in a fixed subspace of integer dimension, here the vertex plane, with the ambient dimension continued. The continuation needs only the area of a sphere and no dynamics of a world with $d \ne 3$. It gives no meaning to a space of non-integer dimension in the register, and none is needed, because it is used only to test the counting equation for other solutions.

\section{The Horizon and the Factor of $\pi$}
\label{sec:pi}

\textit{The sphere is the horizon.} With $d = 3$ the sphere of Postulate~2 is a two-sphere of radius $R = L\ell_P/2\pi$. The value of $L$ is not derived here. Placing the ring on a great circle of the cosmological horizon makes $L$ the circumference of the horizon in Planck lengths, Eq.~(\ref{eq:ident}), and the observed $\Lambda$ then gives $L = 6.4\times10^{61}$. With that value $R = R_\Lambda$, and the sphere is the cosmological horizon. The sphere is about the place of the interaction (Sec.~\ref{sec:postulates}), and the cosmological horizon of de Sitter space has that property: an observer at rest has a horizon of its own, a sphere about the observer whose great circles have length $2\pi R_\Lambda$~\cite{gibbonshawking1977}. The sphere of Postulate~2 is therefore at once the sphere of directions at the vertex, at the scale of the horizon, and the horizon of an observer at rest there. That is the sense in which the sphere belongs to the interaction. From here on the paper calls the sphere the horizon, with its $L^2/\pi$ Planck pixels and its entropy $L^2/4\pi$.

\textit{The residual factor of $\pi$.} The $L^2$ rational cells and the $L^2/\pi$ Planck pixels are two ways of counting one sphere, at two resolutions. The rational grid on $S^2$ is $L$ bands of latitude by $L$ sectors of longitude (Sec.~\ref{sec:supply}). It is an equal-area partition of the sphere into $L^2$ cells, and on the horizon each cell has area
\begin{equation}
\frac{4\pi R_\Lambda^2}{L^2} = \frac{4\pi}{L^2}\Big(\frac{L\ell_P}{2\pi}\Big)^2 = \frac{\ell_P^2}{\pi},
\label{eq:cell}
\end{equation}
which is $0.318$ of a Planck pixel. The register's grid over-resolves the horizon by a factor $\pi$ in area, and that accounts for the residual in full. On the grid the result is exact, Eq.~(\ref{eq:three}). The $\pi$ comes from counting in Planck pixels instead, and Sec.~\ref{sec:root} shows that it amounts to $0.008$ of a dimension. In the grid's cells the horizon entropy $L^2/4\pi$ reads $N_{\rm cell}/4\pi$, an observation left open. The relation also runs the other way. A Planck pixel subtends $\sqrt{4\pi/(L^2/\pi)} = 2\pi/L$ on the horizon, exactly one notch of the register's phase, so along the grid's equator its cell and the Planck pixel have the same length, $\ell_P$. Only the relation between an area and a squared circumference, $A = C^2/\pi$, brings in the $\pi$.

This is the distinction that separates Davies' cosmological information bound~\cite{davies2007}, the horizon area in Planck units, $\log_2(L^2/\pi) = 409$, from Palmer's $\log_2 L = 205$: area against circumference.

\section{Interactions and States}
\label{sec:states}

The number $L^3$ also appears in the framework, and it does not compete with the $L^2$ of Eq.~(\ref{eq:demand}). A configuration of several particles in several planes needs of order $L^3$ labels, two magnitudes and one relative angle for a pair (below), and $L^3 = 4\pi N_{\rm pix}(4)$ to the same accuracy of order one with which $L^2 = \pi N_{\rm pix}(3)$. A reader who applied the rule of one momentum per cell to that number would select $d = 4$, and the argument would then depend on a choice of what to count. The rule applies to an alphabet, and $L^3$ is not one. The alphabet is the set from which a single momentum is drawn. A state is a function on the alphabet, an amplitude for each element, and functions are not elements. A state never competes with the alphabet for cells, and it is not what Postulate~2 is about.

The two numbers differ at exactly the orientation of a vertex (Sec.~\ref{sec:vertex}). For one vertex the orientation of its plane is a choice of coordinates, and the alphabet is $L^2$. For two momenta not confined to one vertex, the overall rotation is a choice but the angle between them is not: two magnitudes and one relative angle make $L^3$, which is that number. So $L^2$ and $L^3$ count different objects, the momentum alphabet of a vertex and the configuration space of a pair, and the second is not a rival calculation of the first. Postulate~2 is about the alphabet of one vertex, and it says nothing about a configuration.

A configuration does not exceed the capacity of the horizon. Its $L^3$ labels are $3\log_2 L = 616$ bits, and the log-dimension of the horizon, the logarithm of the dimension of its Hilbert space, is $L^2/4\pi = 3.3\times10^{122}$ nats. What is open is the encoding. The horizon has $L^2$ cells and a configuration has $L^3$ labels, so the rule of one cell per label, which Postulate~2 states for one vertex, cannot extend to configurations. How states of the interior are written on the horizon is the problem of de Sitter holography~\cite{bousso1999,banksfischler2001,parikh2003}. That problem concerns the Hilbert space and not the dimension of space. This paper does not answer it, and no result here depends on it. The alphabet of a vertex does not compete with that capacity either: an alphabet is not a state, since every interaction reuses it, so Bousso's $N$-bound~\cite{bousso2000} is not invoked for Postulate~2.

\section{The Algebraic Reading}
\label{sec:algebra}

The counting of Sec.~\ref{sec:vertex} has an algebraic counterpart. It rests on an axiom that this paper does not assume, and it is not used in the count; it is recorded because it reaches the same number by another route. What it fixes is the dimension of the Euclidean space whose Clifford algebra is the algebra of one qubit, and not the full Dirac structure of $3+1$ dimensions, for which Sec.~\ref{sec:consequences} says what else is needed.

The shift $X$ has eigenvalues $e^{2\pi i\kappa/L}$, roots of unity, so for $L > 2$ the amplitudes are complex and the number field is $\mathbb{C}$. The cyclic group $\mathbb{Z}_L$ has one automorphism that inverts every translation, $j \mapsto -j$, and it is the momentum doublet of Sec.~\ref{sec:register}: the fundamental object of a free particle on the ring is a two-level system over $\mathbb{C}$. Its observable algebra is $M_2(\mathbb{C})$, and
\begin{align}
\mathrm{Cl}(1,0) &= \mathbb{R}\oplus\mathbb{R}, & \mathrm{Cl}(2,0) &= M_2(\mathbb{R}), \nonumber\\
\mathrm{Cl}(3,0) &= M_2(\mathbb{C}), & \mathrm{Cl}(4,0) &= M_2(\mathbb{H}),
\label{eq:clifford}
\end{align}
so $M_2(\mathbb{C})$ is the Clifford algebra of a three-dimensional space and of no other dimension, $\mathrm{Cl}(3,0)$ in Euclidean signature and $\mathrm{Cl}(1,2)$ in the one other three-dimensional signature that yields it~\cite{lounesto2001}. Real amplitudes would give two dimensions and quaternionic ones four, so the number field selects the dimension.

The axiom this reading needs is that the three Clifford generators are spatial directions. The doublet supplies one of them, since $\sigma_z$ labels the two senses of motion along the ring, and Palmer's identification of a measurement setting with a point of the Bloch sphere~\cite{palmer2020,palmer2026c} supplies all three. With that axiom the algebra selects three dimensions with no reference to interactions.

This route and the counting are one structure read twice. Both descend from the same $\mathbb{Z}_L$ and the same doublet, and the two factors of $L$ in Eq.~(\ref{eq:demand}) are the heading and the magnitude. Their agreement on three is a consistency condition the framework had to pass, and it is not two determinations.

The grid is rational in two angles because a qubit's state space is a two-sphere. Palmer's identification of a spin analyzer's direction with a point on that sphere is possible only when the sphere of spatial directions is also a two-sphere, which is $d = 3$, and which he assumes. Here $d = 3$ is counted from the legs of a vertex and is not assumed, so his identification is permitted by the result and is not used to obtain it.

The register's finite rotational structure can be named. The grid of Eq.~(\ref{eq:raqm}) is invariant under rotations about the pole by multiples of $2\pi/L$, and under rotations by $\pi$ about equatorial axes at permitted azimuths, which exchange the two senses of motion (Sec.~\ref{sec:register}). Together they form the dihedral group of order $2L$. The grid is not invariant under $SO(3)$: in three dimensions, three non-coplanar settings cannot all lie on the grid at once, which is Palmer's Impossible Triangle Corollary~\cite{palmer2026b,palmer2026c} in spatial form.

In $3+1$ dimensions the Dirac matrices generate the complexification of $\mathrm{Cl}(3,1)$, which is $M_2(\mathbb{C})\otimes M_2(\mathbb{C})$, and in it chirality and spin are separate qubits. The register supplies chirality, and Sec.~\ref{sec:consequences} names the transverse qubit that supplies spin, where it is assumed and not derived.

\section{Consequences for the Register}
\label{sec:consequences}

With $d = 3$ in hand, five things follow for the description of a particle on the register, which the companion paper gives along one line.

\textit{One ring per direction, and the circle of headings has its own pair.} A momentum along a heading is a mode on the ring of that direction (Sec.~\ref{sec:register}). Each ring carries its own Weyl pair, the shift $X$ and the operator $Z$ that multiplies $|j\rangle$ by $e^{2\pi ij/L}$, with $ZX = e^{2\pi i/L}XZ$, so $\Delta x\,\Delta p \ge \hbar/2$ holds along every direction by the Donoho--Stark argument~\cite{donoho1989,massar2008}, a state on $N_x$ cells occupying at least $L/N_x$ momenta. The heading itself is a cell on a great circle, and that circle carries a Weyl pair too. Its cell is an angle, $2\pi/L$, so the conjugate quantum is $2\pi\hbar/(L\cdot2\pi/L) = \hbar$. The generator of one-notch rotations of the heading, which is the angular momentum about the plane's normal, therefore takes the values $\hbar\nu$ with $\nu\in\mathbb{Z}_L$, and $[\hat\phi,\hat L_z] = i\hbar$ is the same finite Weyl relation as $[\hat x,\hat p] = i\hbar$ on the ring. The quantization of orbital angular momentum in units of $\hbar$ is thus not imported.

Its range provides a check. The values $|\nu| \le L/2$ give a largest orbital angular momentum $\hbar L/2 = 3.4\times10^{27}$~J\,s. The largest a particle can carry is the largest momentum on the ring, at $\kappa = L/2$, times the horizon radius, $(\pi\hbar/\ell_P)(L\ell_P/2\pi) = \hbar L/2$, the same number.

\textit{Isotropy, and what a rotation is.} A fixed lattice picks out axes, and the register is not a fixed lattice. Each mode advances $|\kappa|$ notches of phase per tick along its own ring (Sec.~\ref{sec:register}), so the dispersion is the same in every permitted direction, and the register's prediction of no energy-dependent speed of light extends from one dimension to three. No physical direction is preferred: the pole of the grid is a choice of coordinates (Sec.~\ref{sec:register}), and the normal of any plane can serve as it. What the grid lacks is a continuum of rotation angles. Its symmetry group is the dihedral group of Sec.~\ref{sec:algebra} and not $SO(3)$. That is the discretization the framework is built on and not a lattice defect, the same fact as the Impossible Triangle Corollary~\cite{palmer2020,palmer2026b}. Rotational invariance in the continuum sense is therefore in the same state as Lorentz invariance, asserted by Palmer and deferred~\cite{palmer2026c}. The grid restricts which headings a state can have, not how a mode moves along its heading, so the tests that bound an energy-dependent speed of light or a preferred direction bound nothing here; the register predicts neither.

\textit{Momentum in the plane, position on the ring.} In a plane, momentum has the clean parameterization of Eq.~(\ref{eq:demand}), and position in whole cells does not, because shifts along non-collinear directions are incommensurate apart from Pythagorean pairs. The orientation of a vertex (Sec.~\ref{sec:vertex}) removes only the overall orientation, which is a symmetry of the state space and not a reduction of it.

\textit{Spin.} In one dimension the doublet is chirality. In three dimensions the plane transverse to the heading $\hat{\mathbf{k}}$ exists. Wigner's little group of a massive particle, the rotations that leave its momentum unchanged, is $SU(2)$, and its fundamental representation is one qubit, spin (Ref.~\cite{weinberg1995}, Sec.~2.5). The four components of the Dirac field are $2\times2$: chirality from the register, and spin transverse to the heading from the third dimension. For a massless mode the little group's rotation is the $U(1)$ of helicity, the component of spin along the heading, the two coincide, and the register's doublet is helicity. The transverse qubit's azimuth is a pattern on the great circle perpendicular to $\hat{\mathbf{k}}$, so the register's clause that every qubit is a pattern on a ring holds. Its parameters are those of Eq.~(\ref{eq:raqm}) on that circle: the frequency of spin up along $\hat{\mathbf{k}}$ is $m/L$, and the azimuth of the spin about $\hat{\mathbf{k}}$ is $2\pi n/L$. With $d = 3$ its state can be identified with a direction of space, the identification that Palmer makes; neither this paper nor the companion paper needs it. The transverse qubit itself is assumed here and not derived. The count of Sec.~\ref{sec:vertex} does not use it.

\textit{Capacity.} A state of many qubits is many strings read jointly, cell by cell (Sec.~\ref{sec:postulates}), in blocks. A later string is read on the cells of each block of earlier outcomes, its fraction of 1s there being the conditional probability of its qubit, and two branches of earlier outcomes share a block only if the state gives the qubit the same conditional on both, in fraction and in phase. A block whose conditional lies strictly between 0 and 1 needs two cells, and a generic state gives its last qubit a different conditional on each of the $2^{N-1}$ branches of the first $N-1$, so $2^N \le L$, a capacity $N_{\max} = \lfloor\log_2L\rfloor = 205$ that the companion paper develops. The count is over the blocks of one state on one ring. It says nothing about how correlations are spread over space and needs no field theory, so the dimension leaves it untouched. What changes is a third role of $L$, the dimension of a free particle's orbital Hilbert space, which is $L$ on the ring. In three dimensions the orbital dimension of one free particle is of order $L^3$, whose $\log_2$ is $616$. The ceiling counts outcomes and not that dimension.

\section{Predictions and Falsifiers}
\label{sec:pred}

The result for the dimension is a postdiction: $d = 3$ for an observed $3$, from the three legs of the register's interaction. Its content is the relation between the two numbers, Eq.~(\ref{eq:nk}), and the sharpness of Sec.~\ref{sec:root}.

\textit{Agreement with known interactions.} The register's own vertex has three legs (Sec.~\ref{sec:vertex}). So do the elementary vertices of the Standard Model apart from the scalar quartic: its three-point vertices directly, and its four-point terms fixed by a cubic coupling, the gauge seagull and the four-gluon vertex, as the square. Contact interactions with five or more legs are non-renormalizable in $3+1$ dimensions. A decay is one string in and two out, three legs. A three-body decay or a two-body annihilation is two such vertices with a line between them. What the framework does not construct is the decay itself. The register's dynamics emits and absorbs massless quanta, and decays of matter are left to a later paper.

\textit{Falsifiers.} The result has the falsifiers of its two postulates. Postulate~1 makes $L$ universal, and the count sets the same $L$ on both sides, so the result stands or falls with a universal $L$. No state then spans more than $205$ qubits (Sec.~\ref{sec:consequences}), and a random circuit and its inverse that return past $205$ logical qubits would falsify it. An energy-dependent speed of light at any order would falsify Postulate~1 as well, since its shift moves every mode one cell per tick. Postulate~2 adds two tests of its own. It makes every great circle the same ring, so the speed of light is the same in every direction (Sec.~\ref{sec:consequences}), and a dependence on direction would falsify it. With $L$ from the cosmological constant it makes the sphere the cosmological horizon (Sec.~\ref{sec:pi}), so the ceiling on qubits and the circumference of the horizon are one number. A ceiling found well away from $205$ would leave $L$ universal and break that identification. A vertex with four legs is not a falsifier. Its momenta would need four dimensions, and it would look three-dimensional from outside (Sec.~\ref{sec:vertex}), and its only trace would be a coupling of its own, one that is not the product of cubic couplings. The scalar quartic is the place to look, since the Standard Model's $\lambda$ is an independent coupling as written.

\section{Discussion}
\label{sec:discussion}

The register, Postulate~1, is a statement about a circle, and it says nothing about the sphere on which the circles lie. The second postulate is about that sphere: its great circles are rings, and its cells label the momenta of an elementary interaction. Counting both sides gives the interaction as many dimensions as it has legs. The register's interaction has three, and every observation is made with it, so the space that is observed has three dimensions, whatever else takes place.

The ledger of the argument is as follows, and it is also the list of what the argument does not do. Postulated: the register with a zero mark of phase at every cell (Postulate~1), and the sphere whose cells label the momenta of an elementary interaction (Postulate~2), read as a sphere about the place of the interaction. Assumed: a background with a linear structure at the point of an interaction, in which a direction is fixed by $d-1$ numbers and vectors add; this is less than the locally Euclidean space of Refs.~\cite{muller2013,dakic2016}, and whether the linearity itself can be derived is asked in Ref.~\cite{muller2017}. Taken from quantum mechanics: the Hilbert space, with the Bloch sphere of a qubit. Constructed: the rational grid on a sphere of $d-1$ angles, Palmer's rule applied to each (Sec.~\ref{sec:supply}). Derived here from Postulate~1: Born's rule as a frequency over the string, the $L$ modes of a ring, one ring for each direction of motion, that the orientation of the grid is a choice of coordinates, the conservation of momentum, and the form of the register's least interaction. Also from Postulate~1, with the joint reading of many strings in blocks (Sec.~\ref{sec:consequences}): the capacity $2^N \le L$. Taken from quantum mechanics, with the Dirac Hamiltonian: the momentum doublet, which the count uses only for the two energy branches, and the tick with its flip, which the count does not use; the companion paper develops both along one line. Taken from field theory: that the gauge field of a free choice of phase is the electromagnetic field, which identifies the waves of the offsets with light.

Derived from both postulates: that the register's interaction has three legs, the third carried by the offsets on the sphere, whose closed-path sums are data, are changed by a quantum, and travel as a wave with the momenta of a ring. Also from both: that those legs cannot be collinear; that a vertex with $k$ legs needs $k$ dimensions, so that $d = 3$ for the register's interaction; and that every record is made with that interaction, up to one pair in $L$. It follows that every observation is three-dimensional, a vertex with four legs included. Also derived: that two uncorrelated strings on one ring meet at a cell for exactly one pair in $L$, and for fewer off it; the radius $L\ell_P/2\pi$ of the sphere; and that the count relaxed to real $d$ has one root, within the bound that the relaxation itself imposes. Not derived: the placement of the ring on the horizon, which fixes the value of $L$ from the observed $\Lambda$; the de Sitter horizon, taken from general relativity; the universality of $d$; and space itself, since the background is assumed and the self-consistent $d$ is found. Open: the coupling and the equation of motion of the offsets, the relation of the grid's cells to the horizon entropy, the encoding of many-particle states on the horizon (Sec.~\ref{sec:states}), the four-point contact terms of field theory with their strengths, the scalar quartic among them, the origin of the transverse qubit, time, and signature.

\begin{acknowledgments}
The author used an AI assistant (Anthropic's Claude) for adversarial review, numerical verification, and editing of the manuscript. The physical proposal, the postulates, and the claims are the author's own. This research received no external funding. The author declares no conflicts of interest. No new data were created or analyzed in this study.
\end{acknowledgments}

\begin{thebibliography}{99}
\bibitem{korytko2026c} A. Korytko, One register: a common origin for the granularity of Hilbert space and of space, Zenodo preprint (2026). \mbox{\url{https://doi.org/10.5281/zenodo.22669417}}
\bibitem{palmer2020} T. N. Palmer, Discretization of the Bloch sphere, fractal invariant sets and Bell's theorem, Proc. R. Soc. A \textbf{476}, 20190350 (2020).
\bibitem{palmer2026} T. N. Palmer, Rational quantum mechanics: Testing quantum theory with quantum computers, Proc. Natl. Acad. Sci. USA \textbf{123}, e2523350123 (2026); arXiv:2510.02877.
\bibitem{palmer2026b} T. N. Palmer, Impossible counterfactuals, discrete Hilbert space and Bell's theorem, J. Phys.: Conf. Ser. \textbf{3189}, 012006 (2026); arXiv:2601.14941.
\bibitem{palmer2026c} T. N. Palmer, Solving the mysteries of quantum mechanics: why nature abhors a continuum, arXiv:2602.16382v1 (2026).
\bibitem{ehrenfest1917} P. Ehrenfest, In what way does it become manifest in the fundamental laws of physics that space has three dimensions?, Proc. Amsterdam Acad. \textbf{20}, 200 (1917).
\bibitem{lovelace1971} C. Lovelace, Pomeron form factors and dual Regge cuts, Phys. Lett. B \textbf{34}, 500 (1971).
\bibitem{goddard1972} P. Goddard and C. B. Thorn, Compatibility of the dual Pomeron with unitarity and the absence of ghosts in the dual resonance model, Phys. Lett. B \textbf{40}, 235 (1972).
\bibitem{wolfram2020} S. Wolfram, A class of models with the potential to represent fundamental physics, Complex Systems \textbf{29}, 107 (2020); arXiv:2004.08210.
\bibitem{muller2013} M. P. M\"uller and Ll. Masanes, Three-dimensionality of space and the quantum bit: an information-theoretic approach, New J. Phys. \textbf{15}, 053040 (2013).
\bibitem{masanes2014} Ll. Masanes, M. P. M\"uller, D. P\'erez-Garc\'ia, and R. Augusiak, Entanglement and the three-dimensionality of the Bloch ball, J. Math. Phys. \textbf{55}, 122203 (2014).
\bibitem{dakic2016} B. Daki\'c and \v{C}. Brukner, The classical limit of a physical theory and the dimensionality of space, in \emph{Quantum Theory: Informational Foundations and Foils}, edited by G. Chiribella and R. W. Spekkens (Springer, Dordrecht, 2016), Ch.~8.
\bibitem{dariano2014} G. M. D'Ariano and P. Perinotti, Derivation of the Dirac equation from principles of information processing, Phys. Rev. A \textbf{90}, 062106 (2014).
\bibitem{muller2017} M. P. M\"uller, S. Carrozza, and P. A. H\"ohn, Is the local linearity of space-time inherited from the linearity of probabilities?, J. Phys. A: Math. Theor. \textbf{50}, 054003 (2017).
\bibitem{niven1956} I. Niven, \emph{Irrational Numbers} (Mathematical Association of America, Washington, DC, 1956).
\bibitem{jahnel2010} J. Jahnel, When is the (co)sine of a rational angle equal to a rational number?, arXiv:1006.2938 (2010).
\bibitem{bousso2000} R. Bousso, Positive vacuum energy and the $N$-bound, J. High Energy Phys. \textbf{11}, 038 (2000).
\bibitem{thooft1972} G. 't Hooft and M. Veltman, Regularization and renormalization of gauge fields, Nucl. Phys. B \textbf{44}, 189 (1972).
\bibitem{wilson1974} K. G. Wilson, Confinement of quarks, Phys. Rev. D \textbf{10}, 2445 (1974).
\bibitem{peierls1933} R. Peierls, Zur Theorie des Diamagnetismus von Leitungselektronen, Z. Phys. \textbf{80}, 763 (1933).
\bibitem{kogut1975} J. Kogut and L. Susskind, Hamiltonian formulation of Wilson's lattice gauge theories, Phys. Rev. D \textbf{11}, 395 (1975).
\bibitem{davies2007} P. C. W. Davies, The implications of a cosmological information bound for complexity, quantum information and the nature of physical law, Fluct. Noise Lett. \textbf{7}, C37 (2007).
\bibitem{bousso1999} R. Bousso, Quantum global structure of de Sitter space, Phys. Rev. D \textbf{60}, 063503 (1999).
\bibitem{banksfischler2001} T. Banks and W. Fischler, M-theory observables for cosmological space-times, arXiv:hep-th/0102077 (2001).
\bibitem{parikh2003} M. K. Parikh, I. Savonije, and E. Verlinde, Elliptic de Sitter space: dS/$\mathbb{Z}_2$, Phys. Rev. D \textbf{67}, 064005 (2003).
\bibitem{gibbonshawking1977} G. W. Gibbons and S. W. Hawking, Cosmological event horizons, thermodynamics, and particle creation, Phys. Rev. D \textbf{15}, 2738 (1977).
\bibitem{planck2018} Planck Collaboration, N. Aghanim et al., Planck 2018 results. VI. Cosmological parameters, Astron. Astrophys. \textbf{641}, A6 (2020).
\bibitem{lounesto2001} P. Lounesto, \emph{Clifford Algebras and Spinors}, 2nd ed. (Cambridge University Press, 2001), Ch.~16.
\bibitem{donoho1989} D. L. Donoho and P. B. Stark, Uncertainty principles and signal recovery, SIAM J. Appl. Math. \textbf{49}, 906 (1989).
\bibitem{massar2008} S. Massar and P. Spindel, Uncertainty relation for the discrete Fourier transform, Phys. Rev. Lett. \textbf{100}, 190401 (2008).
\bibitem{weinberg1995} S. Weinberg, \emph{The Quantum Theory of Fields}, Vol.~1 (Cambridge University Press, 1995).
\end{thebibliography}

\end{document}
